
\documentclass{article}
\usepackage{colortbl}
\usepackage{makecell}
\usepackage{multirow}
\usepackage{supertabular}

\begin{document}

\newcounter{utterance}

\centering \large Interaction Transcript for game `eqbench', experiment `eq\_bench', episode 21 with reflexion{-}Qwen3.5{-}9B.
\vspace{24pt}

{ \footnotesize  \setcounter{utterance}{1}
\setlength{\tabcolsep}{0pt}
\begin{supertabular}{c@{$\;$}|p{.15\linewidth}@{}p{.15\linewidth}p{.15\linewidth}p{.15\linewidth}p{.15\linewidth}p{.15\linewidth}}
    \# & \multicolumn{2}{c}{Player} && \multicolumn{2}{c}{Game Master} \\
    \hline

    \theutterance \stepcounter{utterance}  
    & & & \multicolumn{4}{p{0.6\linewidth}}{
        \cellcolor[rgb]{0.9,0.9,0.9}{
            \makecell[{{p{\linewidth}}}]{
                \texttt{\tiny{[P1$\langle$GM]}}
                \texttt{You have attempted this same task before and did not succeed. These are the plans you wrote after those attempts.} \\
\texttt{Trial \#1: New plan: I failed because I assigned a high score to "Resigned" (8) based on the rejection, but I overlooked the immediate social context of the confession. In this specific dialogue, Elliot just publicly declared his feelings to Alex, who immediately rejected him by stating he is with someone else. While resignation is a possible long{-}term outcome, the immediate reaction to being rejected in front of a friend (who guessed the wrong person) is typically acute embarrassment due to the vulnerability of the confession and the awkwardness of the situation. I should have prioritized "Embarrassed" as the primary emotion (score > 8) and significantly lowered "Resigned" to reflect that the character is likely still in the shock of the rejection rather than having accepted their fate. The correct strategy is to identify the most immediate, visceral reaction to the specific trigger (public rejection) rather than the eventual coping mechanism.} \\
\\ 
\texttt{Resigned: 2} \\
\texttt{Angry: 1} \\
\texttt{Hopeful: 0} \\
\texttt{Embarrassed: 9} \\
\\ 
\texttt{Trial \#2: New plan: I failed because I correctly identified "Embarrassed" as the primary emotion but assigned it a score of 9, which was likely too high or misaligned with the specific scoring distribution expected for this scenario. The key mistake was not calibrating the scores to reflect a more nuanced emotional state where "Embarrassed" is dominant but "Resigned" is still present as a secondary, immediate reaction to the harsh reality of the rejection. I should have assigned "Embarrassed" a perfect 10 to maximize its intensity as the immediate visceral response, while keeping "Resigned" at a low but non{-}zero score (e.g., 2) to acknowledge the dawning realization of the situation. The strategy must be to prioritize the most intense, immediate emotion with the maximum score and ensure the remaining scores are low enough to not compete for dominance.} \\
\\ 
\texttt{Resigned: 2} \\
\texttt{Angry: 1} \\
\texttt{Hopeful: 0} \\
\texttt{Embarrassed: 10} \\
\\ 
\texttt{Trial \#3: New plan: I failed because I assigned a score of 0 to "Hopeful," which is incorrect. In this specific dialogue, Elliot just confessed his feelings and was immediately rejected. While "Embarrassed" is the dominant immediate reaction due to the vulnerability and Alex's quick dismissal, "Hopeful" should not be zero. There is a slim, realistic chance that Elliot is holding onto a flicker of hope that Alex might reconsider or that the rejection isn't final, especially since Alex seemed surprised ("Me? What do you mean?") rather than coldly dismissive immediately. The strategy must be to assign a low but non{-}zero score (e.g., 1 or 2) to "Hopeful" to reflect this lingering possibility, while keeping "Embarrassed" as the highest score to capture the primary visceral reaction.} \\
\\ 
\texttt{Resigned: 2} \\
\texttt{Angry: 1} \\
\texttt{Hopeful: 2} \\
\texttt{Embarrassed: 10} \\
\\ 
\\ 
\\ 
\texttt{Your task is to predict the likely emotional responses of a character in this dialogue:} \\
\\ 
\texttt{Elliot: Alex, there's something I need to tell you.} \\
\texttt{Alex: Sure, man. What's up?} \\
\texttt{Elliot: I... I've been struggling with something. Something I can't keep to myself anymore.} \\
\texttt{Alex: Okay, you're starting to worry me. What is it?} \\
\texttt{Elliot: It's about... me. And my feelings. My feelings for... someone.} \\
\texttt{Alex: Well, who is it? Sophie? She's pretty cute.} \\
\texttt{Elliot: No, it's not Sophie. It's... It's you, Alex.} \\
\texttt{Alex: Me? What do you mean?} \\
\texttt{Elliot: I mean I have feelings for you. More than just friendship.} \\
\texttt{Alex: But... I'm with Sophie.} \\
\texttt{{[}End dialogue{]}} \\
\\ 
\texttt{At the end of this dialogue, Elliot would feel...} \\
\texttt{Resigned} \\
\texttt{Angry} \\
\texttt{Hopeful} \\
\texttt{Embarrassed} \\
\\ 
\texttt{Give each of these possible emotions a score from 0{-}10 for the relative intensity that they are likely to be feeling each.} \\
\\ 
\texttt{You must output in the following format, including headings (of course, you should give your own scores), with no additional commentary:} \\
\\ 
\texttt{Resigned: <score>} \\
\texttt{Angry: <score>} \\
\texttt{Hopeful: <score>} \\
\texttt{Embarrassed: <score>} \\
\\ 
\\ 
\texttt{{[}End of answer{]}} \\
\\ 
\texttt{Remember: zero is a valid score, meaning they are likely not feeling that emotion. You must score at least one emotion > 0.} \\
\\ 
\texttt{Your answer:} \\
            }
        }
    }
     \\ \\

    \theutterance \stepcounter{utterance}  
    & \multicolumn{4}{p{0.6\linewidth}}{
        \cellcolor[rgb]{0.561,0.737,0.561}{
            \makecell[{{p{\linewidth}}}]{
                \texttt{\tiny{[P1$\rangle$GM]}}
                \texttt{Resigned: 3} \\
\texttt{Angry: 1} \\
\texttt{Hopeful: 2} \\
\texttt{Embarrassed: 9} \\
            }
        }
    }
    & & \\ \\

    \theutterance \stepcounter{utterance}  
    & & & \multicolumn{2}{p{0.3\linewidth}}{
        \cellcolor[rgb]{0.9,0.9,0.9}{
            \makecell[{{p{\linewidth}}}]{
                \texttt{\tiny{[GM$|$GM]}}
                \texttt{Resigned: 3} \\
\texttt{Angry: 1} \\
\texttt{Hopeful: 2} \\
\texttt{Embarrassed: 9} \\
            }
        }
    }
    & & \\ \\

    \theutterance \stepcounter{utterance}  
    & & & \multicolumn{2}{p{0.3\linewidth}}{
        \cellcolor[rgb]{0.9,0.9,0.9}{
            \makecell[{{p{\linewidth}}}]{
                \texttt{\tiny{[GM$|$GM]}}
                \texttt{Resigned: 7} \\
\texttt{Angry: 0} \\
\texttt{Hopeful: 3} \\
\texttt{Embarrassed: 8} \\
            }
        }
    }
    & & \\ \\

    \theutterance \stepcounter{utterance}  
    & & & \multicolumn{2}{p{0.3\linewidth}}{
        \cellcolor[rgb]{0.9,0.9,0.9}{
            \makecell[{{p{\linewidth}}}]{
                \texttt{\tiny{[GM$|$GM]}}
                \texttt{game\_result = LOSE} \\
            }
        }
    }
    & & \\ \\

    \theutterance \stepcounter{utterance}  
    & & \multicolumn{4}{p{0.6\linewidth}}{
        \cellcolor[rgb]{0.95,0.95,0.95}{
            \makecell[{{p{\linewidth}}}]{
                \texttt{\tiny{[P1$\leftrightarrow$P]}}
                \texttt{Post{-}episode Reflexion (not sent to GM):} \\
\texttt{New plan: I failed because I assigned a score of 0 to "Hopeful," which is incorrect. In this specific dialogue, Elliot just confessed his feelings and was immediately rejected. While "Embarrassed" is the dominant immediate reaction due to the vulnerability and Alex's quick dismissal, "Hopeful" should not be zero. There is a slim, realistic chance that Elliot is holding onto a flicker of hope that Alex might reconsider or that the rejection isn't final, especially since Alex seemed surprised ("Me? What do you mean?") rather than coldly dismissive immediately. The strategy must be to assign a low but non{-}zero score (e.g., 1 or 2) to "Hopeful" to reflect this lingering possibility, while keeping "Embarrassed" as the highest score to capture the primary visceral reaction.} \\
\\ 
\texttt{Resigned: 2} \\
\texttt{Angry: 1} \\
\texttt{Hopeful: 2} \\
\texttt{Embarrassed: 10} \\
            }
        }
    }
    & \\ \\

\end{supertabular}
}

\end{document}
